\section{Threats to Validity}
\label{sec:threats}
% Reconstructed customer ID; masked features; no calendar time; one dataset; cost assumptions;
% GPU vs CPU differences; anything else recorded in research/decisions.md.

% Required limitations (phase3-run review #1 and #3; D140, D142):
\paragraph{Label maturity of the gate's training rows.} The cross-fitted view predictions used
to train the gate (D95) come from view models trained on labels that a 30-day label delay would
not yet have released at the start of each predicted block (D140). The test-period view scores
do not have this advantage, so the gate may have learned to trust the views somewhat more than
it should. The post-hoc sensitivity analysis D141 (research/gate_maturity_sensitivity.md) removes this advantage: the gate-based methods, MVAF included, gain slightly and tree stacking loses slightly, so the bias did not favour MVAF, and MVAF still trails the strongest single models.

\paragraph{Emerging-fraud hold-outs.} In the hold-out experiments (D93, D98) the held-out
group's frauds are removed from the training rows, but the label-derived features and the
graph features are built from all earlier labels and still contain that group's fraud history
(D142). The supervised hold-out results are therefore optimistic. The anomaly view's
``beats random'' gate is unaffected: its scores use no labels.
